\begin{abstract}
驾驶日志只覆盖场景的一部分，仿真中的新相机位姿和资产状态却会暴露未见表面。补全这些表面时，需要协调世界空间中的背景、规范空间中的车辆与关节人体，并保持日志记录的外观。我们提出 DriveCompose，通过联合补全分类规范查询，构建可组合的高斯驾驶世界。共享恢复前端提供几何、身份和状态证据；局部查询组保留各类内容的空间结构，全局记忆在背景与资产之间交换上下文。条件流填充未知潜码，每步恢复已见潜码，并抑制已知自由空间。固定的逐查询解码器保持已见高斯参数，独立的渲染检查则识别新增内容造成的退化。候选通过检查后，背景与规范资产保持固定，由相机位姿、刚体变换和骨架姿态驱动不同轨迹的渲染，各条轨迹共享同一次补全结果。该方法将生成过程放在离线世界构建阶段，并通过状态驱动装配明确表达后续场景变化。
\end{abstract}
